\section{Experiments and Analysis}

This section specifies protocols, table skeletons, figure placeholders, and
conditional interpretation plans. No ranking or numerical conclusion is asserted
before verification; every result cell is an unresolved ledger placeholder, and
each analysis states the test it will run and the reading that would support or
fail its hypothesis. Table~\ref{tab:hypotheses-map} states, for each claim, its
atomic falsification gate: estimand, comparator, audit population, decision rule,
and the claim-narrowing action on failure.

\subsection{Setup}
We plan to evaluate three backbones---a primary open model, a larger open model as
a scale-trend point, and a closed model as a generalization row---on a unified
memory-system codebase for baseline fairness. \aca{} uses a frozen embedding
encoder with a two-layer head. Evaluation spans four benchmark categories:
streaming task streams (dialogue and agentic), a web-agent benchmark, a long-term
memory benchmark, and a conversational memory benchmark; the pilot uses two
replayable environments. For fairness, all methods share the backbone, retriever,
and top-$k = 4$ budget, with baselines given an equal token budget for their own
mechanism. We report a two-stage macro-average---streams aggregated to a
benchmark-category score, then equal weight across the four categories---with
task-clustered paired BCa intervals and a simultaneous max-baseline comparison
under Holm correction; seeds default to $\{13, 42, 2026\}$, with the final count
fixed by a pre-registered power analysis. The inferential unit is the
\emph{seed/deployment run}: we use hierarchical/block resampling (block = stream,
pairing = seed) and note that task-clustered BCa alone does not cover run-level
variation; the confirmatory multiplicity family (the C1--C4 gates and the
max-baseline contrast) is predefined under Holm FWER control, with FDR reserved for
exploratory sweeps, and the seed count is computed from the independent-pilot power
analysis (\emph{USER\_APPROVAL\_REQUIRED} sizing inputs). Baselines follow a
fidelity checklist: each is marked a \emph{source-faithful reproduction} or a
\emph{standardized credit-variant wrapper}, and results are reported under exact
method names only for source-faithful reproductions. A subset uses a cross-source
judge, of which $200$ examples are human-checked for agreement
$\tbd{JUDGE_HUMAN_AGREEMENT_KAPPA}$. A single main run costs
$\tbd{MAIN_GPU_HOURS_PER_RUN}$, with the full cost table in the appendix.

\paragraph{Cross-task interference (CTI).}
We adopt the definition
\begin{equation}
\mathrm{CTI} = \mathrm{Acc}_B(\text{B-only bank}) - \mathrm{Acc}_B(\text{A}\cup\text{B mixed bank}),
\end{equation}
the drop in domain-$B$ accuracy when a domain-$A$ bank is mixed in; $\mathrm{CTI}\ge0$
indicates interference. The specific A/B domain pairs are not enumerated in the
source and are locked before running (\emph{USER\_APPROVAL\_REQUIRED}).

\paragraph{Reproducibility protocol.}
Before any confirmatory run we lock, and will report, the following items; each is
currently an open \emph{USER\_APPROVAL\_REQUIRED} decision and none is guessed here:
the S2--S4 stream identities and heterogeneity bins; datasets and versions; the
three backbones; stream construction and task ordering; the bank-size schedule and
warm-up length; the fit/calibration/development/sealed-audit split policy; the judge
protocol; the baseline adapters and fidelity labels; the CTI A/B domain pairs; the
max-baseline contrast statistic; the Holm $\alpha$ and exploratory FDR $q$; and the
per-gate CI/threshold rules of Table~\ref{tab:hypotheses-map}.

\subsection{Main results}
Table~\ref{tab:main-results} orders methods by lineage---no-memory, heuristic
credit, learned credit, orthogonal reasoning, and an intervention reference---and
orders the stream columns S1--S5 by increasing heterogeneity; the identities of
S2--S4 are not specified in the source. We test three predictions against the
eventual numbers: whether the gain grows with stream heterogeneity (from
\tbd{MAIN_GAIN_SINGLE_DOMAIN} on the single-domain stream to
\tbd{MAIN_GAIN_MIXED_STREAM} on the mixed stream); whether learned-credit
baselines are reliably better than heuristic ones; and whether \method{} reaches
\tbd{MERIT_FRACTION_OF_RITFULL_GAIN} of the RIT-Full reference gain at a fraction
of its cost. The headline average gain over the strongest baseline is
\tbd{MAIN_AVG_GAIN_OVER_BEST_BASELINE}. The RIT-Full row is an intervention
reference and is treated as a ceiling only for quantities where that status
follows by construction.

\begin{table*}[t]
\centering
\small
\setlength{\tabcolsep}{2.2pt}
\begin{tabular}{lccccccccc}
\toprule
Method & S1 & S2 & S3 & S4 & S5 & W-A & LME & LCM & AVG \\
\midrule
No-Memory ReAct
& \tbdcell{MAIN_NOMEM_S1} & \tbdcell{MAIN_NOMEM_S2} & \tbdcell{MAIN_NOMEM_S3} & \tbdcell{MAIN_NOMEM_S4} & \tbdcell{MAIN_NOMEM_S5} & \tbdcell{MAIN_NOMEM_WAL} & \tbdcell{MAIN_NOMEM_LME} & \tbdcell{MAIN_NOMEM_LCM} & \tbdcell{MAIN_NOMEM_AVG} \\
Full-History Stuffing
& \tbdcell{MAIN_FULLHIST_S1} & \tbdcell{MAIN_FULLHIST_S2} & \tbdcell{MAIN_FULLHIST_S3} & \tbdcell{MAIN_FULLHIST_S4} & \tbdcell{MAIN_FULLHIST_S5} & \tbdcell{MAIN_FULLHIST_WAL} & \tbdcell{MAIN_FULLHIST_LME} & \tbdcell{MAIN_FULLHIST_LCM} & \tbdcell{MAIN_FULLHIST_AVG} \\
Mem0
& \tbdcell{MAIN_MEM0_S1} & \tbdcell{MAIN_MEM0_S2} & \tbdcell{MAIN_MEM0_S3} & \tbdcell{MAIN_MEM0_S4} & \tbdcell{MAIN_MEM0_S5} & \tbdcell{MAIN_MEM0_WAL} & \tbdcell{MAIN_MEM0_LME} & \tbdcell{MAIN_MEM0_LCM} & \tbdcell{MAIN_MEM0_AVG} \\
A-MEM
& \tbdcell{MAIN_AMEM_S1} & \tbdcell{MAIN_AMEM_S2} & \tbdcell{MAIN_AMEM_S3} & \tbdcell{MAIN_AMEM_S4} & \tbdcell{MAIN_AMEM_S5} & \tbdcell{MAIN_AMEM_WAL} & \tbdcell{MAIN_AMEM_LME} & \tbdcell{MAIN_AMEM_LCM} & \tbdcell{MAIN_AMEM_AVG} \\
AWM
& \tbdcell{MAIN_AWM_S1} & \tbdcell{MAIN_AWM_S2} & \tbdcell{MAIN_AWM_S3} & \tbdcell{MAIN_AWM_S4} & \tbdcell{MAIN_AWM_S5} & \tbdcell{MAIN_AWM_WAL} & \tbdcell{MAIN_AWM_LME} & \tbdcell{MAIN_AWM_LCM} & \tbdcell{MAIN_AWM_AVG} \\
ReasoningBank
& \tbdcell{MAIN_RBANK_S1} & \tbdcell{MAIN_RBANK_S2} & \tbdcell{MAIN_RBANK_S3} & \tbdcell{MAIN_RBANK_S4} & \tbdcell{MAIN_RBANK_S5} & \tbdcell{MAIN_RBANK_WAL} & \tbdcell{MAIN_RBANK_LME} & \tbdcell{MAIN_RBANK_LCM} & \tbdcell{MAIN_RBANK_AVG} \\
SkeMex
& \tbdcell{MAIN_SKEMEX_S1} & \tbdcell{MAIN_SKEMEX_S2} & \tbdcell{MAIN_SKEMEX_S3} & \tbdcell{MAIN_SKEMEX_S4} & \tbdcell{MAIN_SKEMEX_S5} & \tbdcell{MAIN_SKEMEX_WAL} & \tbdcell{MAIN_SKEMEX_LME} & \tbdcell{MAIN_SKEMEX_LCM} & \tbdcell{MAIN_SKEMEX_AVG} \\
Memory-R1
& \tbdcell{MAIN_MEMR1_S1} & \tbdcell{MAIN_MEMR1_S2} & \tbdcell{MAIN_MEMR1_S3} & \tbdcell{MAIN_MEMR1_S4} & \tbdcell{MAIN_MEMR1_S5} & \tbdcell{MAIN_MEMR1_WAL} & \tbdcell{MAIN_MEMR1_LME} & \tbdcell{MAIN_MEMR1_LCM} & \tbdcell{MAIN_MEMR1_AVG} \\
Mem-$\alpha$
& \tbdcell{MAIN_MEMALPHA_S1} & \tbdcell{MAIN_MEMALPHA_S2} & \tbdcell{MAIN_MEMALPHA_S3} & \tbdcell{MAIN_MEMALPHA_S4} & \tbdcell{MAIN_MEMALPHA_S5} & \tbdcell{MAIN_MEMALPHA_WAL} & \tbdcell{MAIN_MEMALPHA_LME} & \tbdcell{MAIN_MEMALPHA_LCM} & \tbdcell{MAIN_MEMALPHA_AVG} \\
MemRL
& \tbdcell{MAIN_MEMRL_S1} & \tbdcell{MAIN_MEMRL_S2} & \tbdcell{MAIN_MEMRL_S3} & \tbdcell{MAIN_MEMRL_S4} & \tbdcell{MAIN_MEMRL_S5} & \tbdcell{MAIN_MEMRL_WAL} & \tbdcell{MAIN_MEMRL_LME} & \tbdcell{MAIN_MEMRL_LCM} & \tbdcell{MAIN_MEMRL_AVG} \\
Reflexion
& \tbdcell{MAIN_REFLEXION_S1} & \tbdcell{MAIN_REFLEXION_S2} & \tbdcell{MAIN_REFLEXION_S3} & \tbdcell{MAIN_REFLEXION_S4} & \tbdcell{MAIN_REFLEXION_S5} & \tbdcell{MAIN_REFLEXION_WAL} & \tbdcell{MAIN_REFLEXION_LME} & \tbdcell{MAIN_REFLEXION_LCM} & \tbdcell{MAIN_REFLEXION_AVG} \\
ReasoningBank+MaTTS
& \tbdcell{MAIN_RBANKMATTS_S1} & \tbdcell{MAIN_RBANKMATTS_S2} & \tbdcell{MAIN_RBANKMATTS_S3} & \tbdcell{MAIN_RBANKMATTS_S4} & \tbdcell{MAIN_RBANKMATTS_S5} & \tbdcell{MAIN_RBANKMATTS_WAL} & \tbdcell{MAIN_RBANKMATTS_LME} & \tbdcell{MAIN_RBANKMATTS_LCM} & \tbdcell{MAIN_RBANKMATTS_AVG} \\
\textbf{MERIT (ours)}
& \tbdcell{MAIN_MERIT_S1} & \tbdcell{MAIN_MERIT_S2} & \tbdcell{MAIN_MERIT_S3} & \tbdcell{MAIN_MERIT_S4} & \tbdcell{MAIN_MERIT_S5} & \tbdcell{MAIN_MERIT_WAL} & \tbdcell{MAIN_MERIT_LME} & \tbdcell{MAIN_MERIT_LCM} & \tbdcell{MAIN_MERIT_AVG} \\
RIT-Full reference
& \tbdcell{MAIN_RITFULL_S1} & \tbdcell{MAIN_RITFULL_S2} & \tbdcell{MAIN_RITFULL_S3} & \tbdcell{MAIN_RITFULL_S4} & \tbdcell{MAIN_RITFULL_S5} & \tbdcell{MAIN_RITFULL_WAL} & \tbdcell{MAIN_RITFULL_LME} & \tbdcell{MAIN_RITFULL_LCM} & \tbdcell{MAIN_RITFULL_AVG} \\
\bottomrule
\end{tabular}
\caption{Main-results skeleton with 14 method rows. Column key: S1--S5 are streaming
streams ordered by increasing heterogeneity (S1 single-domain, S5 mixed; the
identities of S2--S4 remain unspecified in the source); W-A is the web-agent
benchmark, LME the long-term-memory benchmark, LCM the conversational-memory
benchmark, and AVG the two-stage macro-average. All result cells are unresolved.}
\label{tab:main-results}
\end{table*}


Figure~\ref{fig:streaming-efficiency} gives the streaming cumulative-success view
of these results, with the planned \rit{}-budget sweep as an inset (detailed under
Efficiency).

% Ledger IDs:
% [[TBD:MAIN_AVG_GAIN_OVER_BEST_BASELINE]]
% [[TBD:FIG4_LATE_STREAM_GAIN_SHARE]]
\begin{figure}[t]
\centering
\figureplaceholder{1.55in}{%
\textbf{Fig. 4 placeholder: Streaming and efficiency}\\[3pt]
Main panel: $x$ = task index; $y$ = cumulative success.\\
Compare MERIT, strongest verified baselines, and the RIT-Full intervention reference.\\
Optional inset: planned RIT-budget $p$ sweep, shown only after data verification.}
\caption{Planned test of whether performance differences concentrate in the later portion of the stream and whether the RIT budget reaches a Pareto-efficient region.}
\label{fig:streaming-efficiency}
\end{figure}


\subsection{Signal-portability check: reward substitution (secondary analysis)}
The contribution vocabulary of the paper is C1--C4 only. Reward portability, the
zero-shot transfer trend (H3), the Pareto-knee analysis (H4), and the boundary
trends (H5) are reported as \emph{secondary analyses} that probe C1--C4, not as
additional claims; they are kept in the main paper. To test whether the improvement
originates at the signal layer rather than in a
particular pipeline, we substitute $\phihat$ for a manager's original outcome
reward, leaving its optimization otherwise unchanged
(Table~\ref{tab:reward-swap}). We test whether this substitution lifts two
outcome-reward managers, which we treat as complementary rather than merely
correlational, by \tbd{REWARD_SWAP_MEMORY_R1_GAIN} and
\tbd{REWARD_SWAP_MEMALPHA_GAIN} respectively. This is a portability check on one
signal, not a separate method; the full table is a candidate for the appendix.

\begin{table}[t]
\centering
\small
\setlength{\tabcolsep}{3pt}
\begin{tabular}{lll}
\toprule
Manager / reward & AVG & $\Delta$ original \\
\midrule
Memory-R1 / original & \tbdcell{PLUGIN_MEMR1_ORIG_AVG} & -- \\
Memory-R1 / $\widehat{\phi}$ & \tbdcell{PLUGIN_MEMR1_PHI_AVG} & \tbdcell{REWARD_SWAP_MEMORY_R1_GAIN} \\
Mem-$\alpha$ / original & \tbdcell{PLUGIN_MEMALPHA_ORIG_AVG} & -- \\
Mem-$\alpha$ / $\widehat{\phi}$ & \tbdcell{PLUGIN_MEMALPHA_PHI_AVG} & \tbdcell{REWARD_SWAP_MEMALPHA_GAIN} \\
\bottomrule
\end{tabular}
\caption{Reward-substitution skeleton. The table will test signal portability; it does not yet establish that any improvement originates at the signal layer.}
\label{tab:reward-swap}
\end{table}


\subsection{Ablation}
Table~\ref{tab:ablation} removes one component at a time and adds two
alternative-signal controls. We read it as a mechanism-consistent sequence---credit
fidelity (CCC) $\rightarrow$ mechanism and interference change (SR@$20\%$, CTI)
$\rightarrow$ task outcome (AVG)---rather than as isolated cells or a proven
mediation. The pivotal test is A2: the identical \aca{} architecture, features,
update rule, and compute budget, trained on a matched observational/outcome label
instead of the interventional label. If the gain persists under A2, the
contribution is not localized to interventional supervision, and the interpretation
must be revised.

\begin{table*}[t]
\centering
\small
\setlength{\tabcolsep}{4pt}
\begin{tabular}{lccccc}
\toprule
Variant & AVG & $\Delta$AVG & SR@20\% ($t=500$) & CCC & CTI \\
\midrule
Full MERIT & \tbdcell{ABL_FULL_AVG} & 0 & \tbdcell{ABL_FULL_SR20} & \tbdcell{ABL_FULL_CCC} & \tbdcell{ABL_FULL_CTI} \\
A1: no recalibration & \tbdcell{ABL_A1_AVG} & \tbdcell{ABL_A1_DELTA} & \tbdcell{ABL_A1_SR20} & \tbdcell{ABL_A1_CCC} & \tbdcell{ABL_A1_CTI} \\
A2: matched obs.-label control & \tbdcell{ABL_A2_AVG} & \tbdcell{ABL_A2_DELTA} & \tbdcell{ABL_A2_SR20} & \tbdcell{ABL_A2_CCC} & \tbdcell{ABL_A2_CTI} \\
A3: no scope gating & \tbdcell{ABL_A3_AVG} & \tbdcell{ABL_A3_DELTA} & \tbdcell{ABL_A3_SR20} & \tbdcell{ABL_A3_CCC} & \tbdcell{ABL_A3_CTI} \\
A4: no governance & \tbdcell{ABL_A4_AVG} & \tbdcell{ABL_A4_DELTA} & \tbdcell{ABL_A4_SR20} & \tbdcell{ABL_A4_CCC} & \tbdcell{ABL_A4_CTI} \\
A5: no usage features & \tbdcell{ABL_A5_AVG} & \tbdcell{ABL_A5_DELTA} & \tbdcell{ABL_A5_SR20} & \tbdcell{ABL_A5_CCC} & \tbdcell{ABL_A5_CTI} \\
A-judge: LLM self-judge & \tbdcell{ABL_AJUDGE_AVG} & \tbdcell{ABL_AJUDGE_DELTA} & \tbdcell{ABL_AJUDGE_SR20} & \tbdcell{ABL_AJUDGE_CCC} & \tbdcell{ABL_AJUDGE_CTI} \\
A7: group intervention (exploratory) & \tbdcell{ABL_A7_AVG} & \tbdcell{ABL_A7_DELTA} & \tbdcell{ABL_A7_SR20} & \tbdcell{ABL_A7_CCC} & \tbdcell{ABL_A7_CTI} \\
\bottomrule
\end{tabular}
\caption{Ablation skeleton. The table tests credit fidelity, mechanism statistics, interference, and task outcomes without presupposing that any component is necessary. A2 is a capacity-matched control: the identical \aca{} architecture, features, update rule, and compute budget, with only the supervision label swapped from the interventional \rit{} label to a matched observational/outcome label.}
\label{tab:ablation}
\end{table*}

\begin{table*}[t]
\centering
\small
\setlength{\tabcolsep}{3pt}
\renewcommand{\arraystretch}{1.1}
\begin{tabular}{@{}l p{1.32in} p{0.66in} p{0.7in} p{1.9in} p{0.9in}@{}}
\toprule
Gate & Estimand & Comparator & Audit pop. & Decision rule (locked pre-run) & Failure action \\
\midrule
G-C1 & SR@20\% slope over $t$; memory-level CCC$(\coutil,\phiaggrit)$ & baseline banks & pilot audit pool & slope-CI excludes $0$ upward \emph{and} CCC upper-CI below the pilot bound (source: $\mathrm{CCC}>0.5$ falsifies) & narrow C1 to heterogeneous streams \\
G-C2 & event- and memory-level CCC; ACA$-$baseline CCC gain; ECE; cost & correlational baseline & sealed final-audit split & \emph{all} hold: event $\rho\ge0.6$ (lower-CI); memory-level CCC lower-CI $>0$; (ACA$-$baseline) CCC lower-CI $>0$; ECE below its lock; cost $\le$ ceiling & C2 unsupported \\
G-C3a & Full-vs-A4 SR@20\% slope difference & A4 (no governance) & mechanism audit & slope-difference CI excludes $0$, Full flatter (Holm) & drop flattening claim \\
G-C3b & CTI with vs.\ without scope gating & A3 (no scope) & A/B mixed-stream & CTI-reduction CI excludes $0$ (Holm) & drop interference claim \\
G-C3c & AVG lift from $\phihat$-reward substitution & original reward & reward-swap runs & lift lower-CI $>0$ (Holm) & report null portability \\
G-C4 & AVG gain over max baseline \emph{and} token overhead & max-baseline contrast & main runs, macro-AVG & MERIT $>$ max-baseline (Holm, lower-CI $>0$) \emph{and} overhead $\le$ ceiling & report parity \\
G-H3 & zero-shot transfer CCC (secondary) & in-domain \aca{} & transfer runs & transfer-CCC lower-CI above its lock (FDR) & report no transfer \\
G-H4 & Pareto knee on budget--AVG curve (secondary) & budget sweep & efficiency runs & knee exists per the marginal-gain rule & drop Pareto claim \\
G-H5 & directional gain vs.\ bank size / heterogeneity / redundancy (secondary) & within-axis & boundary runs & predicted sign holds (CI, FDR); redundancy is Evidence Gap & scope H5 to passing axes \\
\bottomrule
\end{tabular}
\caption{Atomic falsification gates. C1--C4 are confirmatory (Holm FWER); H3--H5 are secondary/exploratory (FDR). G-C2 and G-C4 are conjunctive. This table is design metadata; it contains no measured results, and every threshold/CI rule marked as a lock is set before any run.}
\label{tab:hypotheses-map}
\end{table*}


\subsection{Mechanism and calibration}
Figure~\ref{fig:mechanism-recovery} reuses the axes and statistical definitions of
the baseline diagnostic (Figure~\ref{fig:diagnostic-baseline}) so that the two are
directly comparable. Panel (a) tests whether \aca{}'s sealed-audit
credit--contribution correlation \tbd{ACA_HELDOUT_CCC_T500} improves on the baseline
\tbd{PILOT_CCC_REASONINGBANK_T500} (H2, gated by G-C2); the source gate for H2 is a
Spearman correlation of at least $0.6$. Panel (b) tests whether the superstition rate
is flattened to \tbd{SR_MERIT_T500}, a reduction of
\tbd{SR_REDUCTION_MERIT_T500_PERCENT} relative to the baseline at $t = 500$; the H1
reading (gated by G-C3a) is the \emph{slope} of this trajectory, compared between the
full method and the no-governance ablation across the deployment checkpoints. Calibration is
assessed with a held-out reliability diagram and expected calibration error
\tbd{ECE_ACA_HELDOUT}, per-domain calibration, and a split-half re-test, with
particular attention to sign instability near $\phicf \approx 0$.

% Ledger IDs:
% [[TBD:ACA_HELDOUT_CCC_T500]]
% [[TBD:PILOT_CCC_REASONINGBANK_T500]]
% [[TBD:SR_MERIT_T500]]
% [[TBD:SR_REDUCTION_MERIT_T500_PERCENT]]
\begin{figure}[t]
\centering
\figureplaceholder{1.55in}{%
\textbf{Fig. 5 placeholder: Mechanism recovery}\\[3pt]
(a) $x$: predicted contribution $\phihat$; $y$: \rit{} label $\phirit$ (finite \rit{} output).\\
(b) $x$: task index $t$ (100--500); $y$: SR@20\%.\\
Compare ACA (MERIT) with the correlational baseline.\\
Axis ranges and statistical definitions must be identical to Fig.~\ref{fig:diagnostic-baseline}.}
\caption{Planned mechanism-closure analysis comparing attribution fidelity and superstition-rate trajectories under the same definitions and axes as the baseline diagnostic.}
\label{fig:mechanism-recovery}
\end{figure}


\subsection{Boundary and scaling (secondary analysis)}
We test whether the gain holds under a larger backbone (from
\tbd{SCALING_GAIN_32B} to \tbd{SCALING_GAIN_235B}, with stronger baselines; the
superstition mechanism is hypothesized to persist) and whether an attributor trained
on the pilot environments transfers zero-shot, measured by the transfer correlation
\tbd{ACA_ZEROSHOT_TRANSFER_CCC_WEBARENA} (H3). For the boundary (H5), we plan three
trend lines: gain versus bank size (the small-bank bin has planned gain
\tbd{H5_GAIN_SMALL_BANK_LT100}), gain versus stream heterogeneity (single-domain
gain \tbd{H5_GAIN_SINGLE_DOMAIN}), and gain versus injected redundancy. The
redundancy trend has no registered result identifier and is reported as an
\textsc{Evidence Gap}, kept exploratory; the hypotheses are that leave-one-out
underestimation grows with redundancy and that the group-intervention variant (A7)
mitigates it.

\subsection{Efficiency}
Figure~\ref{fig:streaming-efficiency} plots cumulative success over the stream and,
as an inset shown only after verification, a planned \rit{}-budget sweep over
$p \in \{1, 2, 5, 10\}\%$; the inset is labeled a planned comparison rather than a
measured curve. We test whether performance differences concentrate in the later,
most-polluted portion of the stream---what share \tbd{FIG4_LATE_STREAM_GAIN_SHARE}
of the average gain arises there---and whether the \rit{} budget saturates near the
source default, placing \method{} at a Pareto-efficient knee (H4).
Table~\ref{tab:efficiency} records relative token cost, success, latency, and
memory overhead against the co-occurrence baseline and the intervention reference.

% Citation ledger IDs used to verify competitor rows:
% [[TBD:CITATION_DATA_SHAPLEY]]
% [[TBD:CITATION_INFLUENCE_FUNCTIONS]]
% [[TBD:CITATION_TRACIN]]
% [[TBD:CITATION_CONTEXTCITE]]
% [[TBD:CITATION_REASONINGBANK]]
% [[TBD:CITATION_SKEMEX]]
% [[TBD:CITATION_MEMORY_R1]]
% [[TBD:CITATION_MEMRL]]
\begin{table*}[t]
\centering
\small
\setlength{\tabcolsep}{3.5pt}
\begin{tabular}{lccccc}
\toprule
Method & Causal signal & O(1) online & Recalibration & Scope & Audit \\
\midrule
Data Shapley & \tbdcell{CITATION_DATA_SHAPLEY} & \tbdcell{CITATION_DATA_SHAPLEY} & \tbdcell{CITATION_DATA_SHAPLEY} & \tbdcell{CITATION_DATA_SHAPLEY} & \tbdcell{CITATION_DATA_SHAPLEY} \\
Influence Functions & \tbdcell{CITATION_INFLUENCE_FUNCTIONS} & \tbdcell{CITATION_INFLUENCE_FUNCTIONS} & \tbdcell{CITATION_INFLUENCE_FUNCTIONS} & \tbdcell{CITATION_INFLUENCE_FUNCTIONS} & \tbdcell{CITATION_INFLUENCE_FUNCTIONS} \\
TracIn & \tbdcell{CITATION_TRACIN} & \tbdcell{CITATION_TRACIN} & \tbdcell{CITATION_TRACIN} & \tbdcell{CITATION_TRACIN} & \tbdcell{CITATION_TRACIN} \\
ContextCite & \tbdcell{CITATION_CONTEXTCITE} & \tbdcell{CITATION_CONTEXTCITE} & \tbdcell{CITATION_CONTEXTCITE} & \tbdcell{CITATION_CONTEXTCITE} & \tbdcell{CITATION_CONTEXTCITE} \\
ReasoningBank & \tbdcell{CITATION_REASONINGBANK} & \tbdcell{CITATION_REASONINGBANK} & \tbdcell{CITATION_REASONINGBANK} & \tbdcell{CITATION_REASONINGBANK} & \tbdcell{CITATION_REASONINGBANK} \\
SkeMex & \tbdcell{CITATION_SKEMEX} & \tbdcell{CITATION_SKEMEX} & \tbdcell{CITATION_SKEMEX} & \tbdcell{CITATION_SKEMEX} & \tbdcell{CITATION_SKEMEX} \\
Memory-R1 & \tbdcell{CITATION_MEMORY_R1} & \tbdcell{CITATION_MEMORY_R1} & \tbdcell{CITATION_MEMORY_R1} & \tbdcell{CITATION_MEMORY_R1} & \tbdcell{CITATION_MEMORY_R1} \\
MemRL & \tbdcell{CITATION_MEMRL} & \tbdcell{CITATION_MEMRL} & \tbdcell{CITATION_MEMRL} & \tbdcell{CITATION_MEMRL} & \tbdcell{CITATION_MEMRL} \\
\textbf{MERIT (ours)} & yes & yes & yes & yes & yes \\
\bottomrule
\end{tabular}
\caption{Method-positioning skeleton. Every competitor cell remains TBD until its cited source is checked item by item; TBD must not be interpreted as lack of capability.}
\label{tab:method-positioning}
\end{table*}

\begin{table*}[t]
\centering
\small
\setlength{\tabcolsep}{4pt}
\begin{tabular}{lccccc}
\toprule
Configuration & Relative tokens & AVG & ACA ms/event & Extra VRAM & RIT overhead \\
\midrule
Co-occurrence baseline & $1\times$ & \tbdcell{EFF_BASELINE_AVG} & -- & -- & 0 \\
MERIT & \tbdcell{MERIT_RELATIVE_TOKEN_COST} & \tbdcell{EFF_MERIT_AVG} & \tbdcell{ACA_SCORING_LATENCY_MS} & \tbdcell{ACA_VRAM_OVERHEAD_GB} & \tbdcell{RIT_TOKEN_OVERHEAD_PERCENT} \\
A1 (static attributor) & \tbdcell{A1_RELATIVE_TOKEN_COST} & \tbdcell{EFF_A1_AVG} & \tbdcell{ACA_SCORING_LATENCY_MS} & \tbdcell{ACA_VRAM_OVERHEAD_GB} & \tbdcell{RIT_TOKEN_OVERHEAD_PERCENT} \\
RIT-Full reference & \tbdcell{RITFULL_RELATIVE_TOKEN_COST} & \tbdcell{EFF_RITFULL_AVG} & -- & -- & \tbdcell{FULL_RIT_COST_MULTIPLIER} \\
\bottomrule
\end{tabular}
\caption{Efficiency skeleton. All non-reference values remain unresolved; no Pareto ranking is assigned before verification.}
\label{tab:efficiency}
\end{table*}


\subsection{Case study}
Figure~\ref{fig:case-study} plans two qualitative timelines. In the success case, a
useful multi-hop memory and a superstitious ``retry the same query'' memory are
tracked over deployment; we test whether the latter's inflated utility is assigned
near-zero \rit{} credit and leads to eviction at task
\tbd{CASE_SUPERSTITION_EVICTION_TASK_INDEX}, while the former's scope is
characterized correctly. In the interaction-boundary case, two
\emph{redundant/substitutable} memories encode overlapping copies of the same
sub-procedure, so that either alone suffices; we test whether individual
leave-one-out is near zero for both---even though their joint value is
positive---so that governance nearly removes both, and whether the exploratory
group-intervention analysis recovers their joint value---an honest display of the
AS3 interaction boundary, with no numerical guarantee. We expect a
share \tbd{RESIDUAL_FAILURE_NO_RELEVANT_MEMORY_FRACTION} of \method{}'s remaining
failures to stem from the bank simply lacking relevant experience, which no credit
signal can repair.

% Ledger ID: [[TBD:CASE_SUPERSTITION_EVICTION_TASK_INDEX]]
\begin{figure}[t]
\centering
\figureplaceholder{1.55in}{%
\textbf{Fig. 6 placeholder: Case study and AS3 boundary}\\[3pt]
Timeline A: useful memory versus superstitious memory, with utility, RIT credit, scope, and governance events.\\
Timeline B: redundant/substitutable-memory case and the exploratory group-intervention analysis.\\
Future source: verified case-study event logs.}
\caption{Planned qualitative analysis of a successful governance case and a redundant/substitutable-memory case that exposes the AS3 interaction boundary.}
\label{fig:case-study}
\end{figure}

